\chapter{Réalisation et Mise en Œuvre}

\onehalfspacing

Ce chapitre décrit la concrétisation des choix de conception présentés au chapitre précédent. Il détaille l'environnement de développement, l'implémentation des modules backend et frontend, les interfaces utilisateur réalisées, ainsi que les résultats des tests de validation. L'objectif est de démontrer que la solution développée répond aux exigences fonctionnelles et non-fonctionnelles identifiées en début de projet.

\section{Environnement de développement}

\subsection{Outils utilisés}

\begin{table}[ht]
\centering
\caption{Outils de développement}
\label{tab:outils}
\begin{tabularx}{\textwidth}{|l|l|X|}
\hline
\textbf{Outil} & \textbf{Version} & \textbf{Usage} \\
\hline
Python & 3.11 & Langage backend \\
\hline
Django & 4.2 & Framework web backend \\
\hline
Django REST Framework & 3.15 & API REST \\
\hline
React.js & 18 & Framework frontend \\
\hline
Vite & 5 & Bundler frontend \\
\hline
TailwindCSS & 3 & Styles utilitaires \\
\hline
PostgreSQL & 15 & Base de données \\
\hline
Docker \& Compose & 24 & Conteneurisation \\
\hline
Nginx & 1.25 & Proxy inverse, TLS \\
\hline
GitHub Actions & -- & Pipeline CI (lint + tests) \\
\hline
VS Code & -- & Éditeur de code \\
\hline
Postman & -- & Tests manuels de l'API \\
\hline
\end{tabularx}
\end{table}

\subsection{Structure du projet}

Le projet est organisé en deux répertoires principaux (\texttt{backend/} et \texttt{frontend/}), complétés par la configuration Docker et Nginx. Le backend Django suit une architecture modulaire : chaque domaine fonctionnel est une application Django distincte avec ses propres modèles, vues, sérialiseurs et URLs.

\begin{lstlisting}[style=PythonStyle, language={}, caption={Structure du projet}]
appels-offres/
|-- docker-compose.yml
|-- nginx/
|   |-- nginx.conf
|   `-- certs/         # Certificat SSL auto-signé (démonstration)
|-- backend/
|   |-- config/        # Settings, urls.py racine, wsgi.py
|   `-- apps/
|       |-- authentication/   # User, JWT, RBAC
|       |-- offers/           # OfferCall, machine à états
|       |-- candidates/       # Candidate, Submission
|       |-- documents/        # Upload, AES-256-GCM
|       |-- notifications/    # Notification, check_deadlines
|       |-- reports/          # Export PDF (reportlab)
|       |-- audit/            # AuditLog, signaux Django
|       `-- portal/           # Endpoints publics (AllowAny)
|-- frontend/
|   `-- src/
|       |-- contexts/         # AuthContext
|       |-- pages/            # Login, Dashboard, offres, portail...
|       |-- services/         # Wrappers API par domaine
|       `-- hooks/            # useAuth, useNotifications
`-- .github/workflows/ci.yml  # GitHub Actions
\end{lstlisting}

\subsection{Pipeline d'intégration continue}

À chaque push sur la branche \texttt{main}, GitHub Actions exécute automatiquement :

\begin{enumerate}
    \item \textbf{Lint} : \texttt{flake8} vérifie le style et la qualité du code Python (longueur des lignes, imports inutilisés, variables non utilisées).
    \item \textbf{Tests} : \texttt{pytest-django} exécute les tests d'intégration avec une base de données PostgreSQL dédiée au pipeline.
\end{enumerate}

\begin{figure}[ht]
\centering
\fbox{%
  \begin{minipage}{0.85\textwidth}
    \vspace{2cm}
    \centering
    \textit{[À INSÉRER — Pipeline CI GitHub Actions — capture d'écran]}
    \vspace{2cm}
  \end{minipage}%
}
\caption{Pipeline CI GitHub Actions (lint + tests)}
\label{fig:ci}
\end{figure}

\section{Implémentation du backend}

\subsection{Authentification et RBAC}

Le modèle \texttt{User} étend \texttt{AbstractBaseUser} de Django pour définir le champ \texttt{email} comme identifiant principal (et non le nom d'utilisateur). Un champ \texttt{role} de type \texttt{CharField} avec des choix énumérés (\texttt{ADMIN}, \texttt{MANAGER}, \texttt{ANALYST}, \texttt{VIEWER}) définit le rôle de chaque utilisateur.

\begin{lstlisting}[style=PythonStyle, language=Python, caption={Modèle User}]
class User(AbstractBaseUser, PermissionsMixin):
    id = models.UUIDField(primary_key=True, default=uuid.uuid4)
    email = models.EmailField(unique=True)
    full_name = models.CharField(max_length=150)
    role = models.CharField(max_length=10, choices=ROLE_CHOICES, default='VIEWER')
    department = models.CharField(max_length=100, blank=True)
    is_active = models.BooleanField(default=True)
    created_at = models.DateTimeField(auto_now_add=True)
\end{lstlisting}

Les permissions DRF sont définies comme des classes réutilisables : \texttt{IsManagerOrAdmin} (création et modification d'offres), \texttt{IsAnalystOrAbove} (lecture de toutes les offres, évaluation), \texttt{IsAdminRole} (gestion des utilisateurs, journal d'audit). Ces classes sont appliquées au niveau de chaque vue et centralisées dans \texttt{apps/authentication/permissions.py}.

L'endpoint de connexion génère un access token (durée 15 minutes) retourné dans le corps de la réponse JSON, et un refresh token (durée 7 jours) placé dans un cookie sécurisé. La bibliothèque \texttt{djangorestframework-simplejwt} gère la génération, la validation et la révocation des tokens.

\subsection{Gestion des offres et machine à états}

Le modèle \texttt{OfferCall} implémente une machine à états explicite. Les transitions autorisées sont définies dans un dictionnaire Python qui associe à chaque statut la liste des statuts atteignables.

\begin{lstlisting}[style=PythonStyle, language=Python, caption={Machine à états des offres}]
ALLOWED_TRANSITIONS = {
    'DRAFT':      ['PUBLISHED', 'CANCELLED'],
    'PUBLISHED':  ['EVALUATION', 'CANCELLED'],
    'EVALUATION': ['AWARDED', 'CANCELLED'],
    'AWARDED':    [],
    'CANCELLED':  [],
}

def can_transition_to(self, new_status):
    return new_status in ALLOWED_TRANSITIONS.get(self.status, [])
\end{lstlisting}

Chaque transition est un endpoint POST dédié (\texttt{/publish/}, \texttt{/close/}, \texttt{/award/}, \texttt{/cancel/}). Cette conception interdit toute modification directe du champ \texttt{status} via PATCH, garantissant que la logique de validation est toujours exécutée.

\subsection{Documents et chiffrement AES-256-GCM}

L'upload de documents s'effectue en multipart/form-data. Avant le stockage, le fichier passe par une validation (type MIME dans la liste blanche PDF/DOCX/XLSX, taille inférieure à 10 Mo), puis par le module de chiffrement si le flag \texttt{is\_sensitive} est activé.

L'algorithme AES-256-GCM garantit à la fois la confidentialité (les données sont illisibles sans la clé) et l'intégrité (le tag d'authentification détecte toute modification). Pour chaque fichier chiffré, un nonce (nombre utilisé une seule fois) de 12 octets est généré aléatoirement, ce qui garantit que deux fichiers identiques produisent des chiffrés différents.

\begin{lstlisting}[style=PythonStyle, language=Python, caption={Chiffrement AES-256-GCM}]
def encrypt_file(data: bytes, key: bytes) -> bytes:
    nonce = os.urandom(12)
    aesgcm = AESGCM(key)
    ciphertext = aesgcm.encrypt(nonce, data, None)
    return nonce + ciphertext  # nonce (12B) || ciphertext || tag (16B)
\end{lstlisting}

Le déchiffrement est transparent pour le client : lors du téléchargement, Django détecte \texttt{is\_encrypted=True}, lit les 12 premiers octets comme nonce, déchiffre le reste, et retourne le fichier en clair dans la réponse HTTP. La clé ne transite jamais en dehors du processus Django.

\subsection{Portail public}

L'application \texttt{portal} est isolée : elle ne partage aucune classe de permission ni aucun sérialiseur avec les applications internes. Elle importe uniquement les modèles (\texttt{OfferCall}, \texttt{Candidate}, \texttt{Submission}, \texttt{Document}).

La vue \texttt{PublicSubmitView} implémente trois protections :

\begin{enumerate}
    \item \textbf{Rate limiting} : décorateur \texttt{@ratelimit(key='ip', rate='5/h', block=True)} de la bibliothèque \texttt{django-ratelimit}. Un script automatisé ne peut pas inonder la base de données. (Dans le contexte de démonstration, ce seuil est volontairement bas ; en production, il serait augmenté.)
    \item \textbf{Vérification de deadline} : \texttt{timezone.now() > offer.submission\_deadline} retourne HTTP 400 si la date limite est dépassée.
    \item \textbf{Unicité} : une société ne peut pas soumettre deux fois une proposition pour la même offre (identifiée par email).
\end{enumerate}

\subsection{Audit et notifications}

Le journal d'audit est alimenté par des appels explicites à \texttt{AuditLog.objects.create()} dans chaque vue lors des opérations d'écriture. Chaque entrée enregistre : l'utilisateur (nullable pour le portail), l'action, le type de ressource, son identifiant, les modifications au format JSON et l'adresse IP.

Le scheduler s'exécute dans un conteneur Docker dédié en boucle :

\begin{lstlisting}[style=PythonStyle, language=bash, caption={Boucle du scheduler}]
while true; do
    python manage.py check_deadlines && sleep 3600 || sleep 60;
done
\end{lstlisting}

La commande \texttt{check\_deadlines} interroge la base de données pour toutes les offres dont la deadline approche (dans les 24 heures) et crée les notifications correspondantes pour les responsables assignés.

\section{Implémentation du frontend}

\subsection{Gestion de l'authentification}

Le contexte \texttt{AuthContext} de React centralise l'état d'authentification : l'access token est stocké dans une variable d'état (mémoire), jamais dans \texttt{localStorage} ou \texttt{sessionStorage}. Au montage de l'application, un appel silencieux à \texttt{/api/auth/refresh/} restaure le token depuis le cookie httpOnly si la session est encore valide.

Un intercepteur Axios ajoute automatiquement l'en-tête \texttt{Authorization: Bearer <token>} à chaque requête, et gère le rafraîchissement transparent en cas de réponse 401 (token expiré).

\subsection{Routage protégé par rôle}

Les routes React sont enveloppées dans des composants de garde : \texttt{PrivateRoute} (authentification requise) et \texttt{RoleGuard} (rôle minimum requis). Un utilisateur VIEWER qui tente d'accéder à la page de création d'offre est redirigé vers le tableau de bord.

\subsection{Portail public}

Le portail utilise une instance Axios séparée (\texttt{portalApi.js}) sans en-tête \texttt{Authorization}. Les pages \texttt{/portal} et \texttt{/portal/:id} sont accessibles sans connexion et ne partagent aucun état avec l'application interne.

\section{Interfaces et fonctionnalités}

\subsection{Tableau de bord}

\begin{figure}[ht]
\centering
\fbox{%
  \begin{minipage}{0.85\textwidth}
    \vspace{2cm}
    \centering
    \textit{[À INSÉRER — Tableau de bord principal]}
    \vspace{2cm}
  \end{minipage}%
}
\caption{Tableau de bord principal}
\label{fig:dashboard}
\end{figure}

\subsection{Liste et filtrage des offres}

\begin{figure}[ht]
\centering
\fbox{%
  \begin{minipage}{0.85\textwidth}
    \vspace{2cm}
    \centering
    \textit{[À INSÉRER — Liste des offres avec filtres]}
    \vspace{2cm}
  \end{minipage}%
}
\caption{Liste des offres avec filtres}
\label{fig:liste_offres}
\end{figure}

La page de liste affiche les offres sous forme de tableau paginé avec filtres par statut, catégorie et recherche textuelle. Des badges colorés signalent l'état de chaque offre (vert : publiée, orange : en évaluation, rouge : annulée).

\subsection{Détail d'une offre}

\begin{figure}[ht]
\centering
\fbox{%
  \begin{minipage}{0.85\textwidth}
    \vspace{2cm}
    \centering
    \textit{[À INSÉRER — Détail d'une offre avec onglets]}
    \vspace{2cm}
  \end{minipage}%
}
\caption{Détail d'une offre avec onglets}
\label{fig:detail_offre}
\end{figure}

La page de détail est organisée en onglets : informations générales, soumissions des candidats (avec scores), et documents associés (avec indicateur de chiffrement AES-256).

\subsection{Formulaire de création d'offre}

\begin{figure}[ht]
\centering
\fbox{%
  \begin{minipage}{0.85\textwidth}
    \vspace{2cm}
    \centering
    \textit{[À INSÉRER — Formulaire de création d'offre]}
    \vspace{2cm}
  \end{minipage}%
}
\caption{Formulaire de création d'offre}
\label{fig:creation_offre}
\end{figure}

\subsection{Évaluation d'une soumission}

\begin{figure}[ht]
\centering
\fbox{%
  \begin{minipage}{0.85\textwidth}
    \vspace{2cm}
    \centering
    \textit{[À INSÉRER — Évaluation et scoring d'une soumission]}
    \vspace{2cm}
  \end{minipage}%
}
\caption{Évaluation et scoring d'une soumission}
\label{fig:evaluation}
\end{figure}

\subsection{Portail public}

\begin{figure}[ht]
\centering
\fbox{%
  \begin{minipage}{0.85\textwidth}
    \vspace{2cm}
    \centering
    \textit{[À INSÉRER — Portail public : liste des offres]}
    \vspace{2cm}
  \end{minipage}%
}
\caption{Portail public — liste des offres}
\label{fig:portail}
\end{figure}

\begin{figure}[ht]
\centering
\fbox{%
  \begin{minipage}{0.85\textwidth}
    \vspace{2cm}
    \centering
    \textit{[À INSÉRER — Formulaire de soumission publique]}
    \vspace{2cm}
  \end{minipage}%
}
\caption{Formulaire de soumission publique}
\label{fig:portail_submit}
\end{figure}

\subsection{Journal d'audit}

\begin{figure}[ht]
\centering
\fbox{%
  \begin{minipage}{0.85\textwidth}
    \vspace{2cm}
    \centering
    \textit{[À INSÉRER — Journal d'audit]}
    \vspace{2cm}
  \end{minipage}%
}
\caption{Journal d'audit}
\label{fig:audit}
\end{figure}

\subsection{Rapport PDF}

\begin{figure}[ht]
\centering
\fbox{%
  \begin{minipage}{0.85\textwidth}
    \vspace{2cm}
    \centering
    \textit{[À INSÉRER — Exemple de rapport PDF généré]}
    \vspace{2cm}
  \end{minipage}%
}
\caption{Exemple de rapport PDF généré}
\label{fig:rapport_pdf}
\end{figure}

\section{Tests et validation}

\subsection{Tests automatisés}

Les tests d'intégration couvrent les scénarios critiques du système :

\begin{table}[ht]
\centering
\caption{Couverture des tests automatisés}
\label{tab:tests}
\begin{tabularx}{\textwidth}{|l|X|}
\hline
\textbf{Module} & \textbf{Cas testés} \\
\hline
Authentification & Connexion valide, identifiants incorrects, refresh token, déconnexion \\
\hline
RBAC & Accès autorisé ou refusé par rôle sur les endpoints critiques \\
\hline
Offres & Création, transitions d'état autorisées et refusées, filtres \\
\hline
Portail & Soumission valide, deadline dépassée, doublon, fichier invalide \\
\hline
Documents & Upload, téléchargement, détection du chiffrement \\
\hline
\end{tabularx}
\end{table}

Les tests utilisent des fixtures par rôle (\texttt{admin\_user}, \texttt{manager\_user}, \texttt{analyst\_user}, \texttt{viewer\_user}) et une base de données PostgreSQL dédiée. Chaque test est isolé dans une transaction annulée après exécution. L'ensemble des tests passe au vert sur le pipeline CI (figure~\ref{fig:ci}).

\subsection{Tests manuels}

Un jeu de données de démonstration a été créé pour valider manuellement les parcours utilisateur : utilisateurs de chaque rôle, offres à différents stades, candidats et soumissions. Ce jeu a permis de tester le parcours complet : création d'offre → publication → soumission via le portail → évaluation → attribution → export PDF → consultation du journal d'audit.

\section{Conclusion}

Ce chapitre a présenté la mise en œuvre complète du système, de l'environnement de développement jusqu'à la validation par les tests. L'implémentation backend a démontré la robustesse de la machine à états des offres, la transparence du chiffrement AES-256-GCM et l'isolation effective du portail public. Côté frontend, la gestion sécurisée des tokens JWT et le routage par rôle garantissent une expérience utilisateur cohérente avec les exigences de sécurité. Les tests automatisés et le pipeline CI confirment la fiabilité de l'ensemble, et les tests manuels valident les parcours fonctionnels de bout en bout.

\newpage
